\chapter{2016年山东卷（文）}

\section{单选题}

\begin{problem}
设全集 \(U=\{1,2,3,4,5,6\}\)，\(A=\{1,3,5\}\)，\(B=\{3,4,5\}\)，则 \(\complement_U\left(A\cup B\right)=\)（\quad）

\choices
    {\(\{2,6\}\)}
    {\(\{3,6\}\)}
    {\(\{1,3,4,5\}\)}
    {\(\{1,2,4,6\}\)}
\end{problem}

\begin{problem}
若复数 \(z=\dfrac{2}{1-\i}\)，其中 \(\i\) 为虚数单位，则 \(z=\)（\quad）

\choices
    {\(1+\i\)}
    {\(1-\i\)}
    {\(-1+\i\)}
    {\(-1-\i\)}
\end{problem}

\begin{problem}
某高校调查了 \(200\) 名学生每周的自习时间（单位：小时），制成了如图所示的频率分布直方图，其中自习时间的范围是 \([17.5,30]\)，样本数据分组为 \([17.5,20)\)，\([20,22.5)\)，\([22.5,25)\)，\([25,27.5)\)，\([27.5,30]\). 根据直方图，这 \(200\) 名学生中每周的自习时间不少于 \(22.5\) 小时的人数是（\quad）

\bitmapfigure[max width=5.8cm]{img/2016/shandong_liberal/q3_fig1.png}

\choices
    {\(56\)}
    {\(60\)}
    {\(120\)}
    {\(140\)}
\end{problem}

\begin{problem}
若变量 \(x\)，\(y\) 满足
\(\begin{cases}
x+y\leq 2,\\
2x-3y\leq 9,\\
x\geq 0,
\end{cases}\)
则 \(x^2+y^2\) 的最大值是（\quad）

\choices
    {\(4\)}
    {\(9\)}
    {\(10\)}
    {\(12\)}
\end{problem}

\begin{problem}
一个由半球和四棱锥组成的几何体，其三视图如图所示. 则该几何体的体积为（\quad）

\examfiguregroup[float=inline]{img/2016/shandong_liberal/q5_fig1.png}{%
    \bitmapinclude[max width=0.48\linewidth]{img/2016/shandong_liberal/q5_fig1.png}\hfill
    \bitmapinclude[max width=0.30\linewidth]{img/2016/shandong_liberal/q5_fig2.png}%
}

\choices
    {\(\dfrac13+\dfrac{2\pi}{3}\)}
    {\(\dfrac13+\dfrac{\sqrt2}{3}\pi\)}
    {\(\dfrac13+\dfrac{\sqrt2}{6}\pi\)}
    {\(1+\dfrac{\sqrt2}{6}\pi\)}
\end{problem}

\begin{problem}
已知直线 \(a\)，\(b\) 分别在两个不同的平面 \(\alpha\)，\(\beta\) 内. 则“直线 \(a\) 和直线 \(b\) 相交”是“平面 \(\alpha\) 和平面 \(\beta\) 相交”的（\quad）

\choices
    {充分不必要条件}
    {必要不充分条件}
    {充要条件}
    {既不充分也不必要条件}
\end{problem}

\begin{problem}
已知圆 \(M:x^2+y^2-2ay=0\)（\(a>0\)）截直线 \(x+y=0\) 所得线段的长度是 \(2\sqrt2\)，则圆 \(M\) 与圆 \(N:(x-1)^2+(y-1)^2=1\) 的位置关系是（\quad）

\choices
    {内切}
    {相交}
    {外切}
    {相离}
\end{problem}

\begin{problem}
\(\triangle ABC\) 中，角 \(A\)，\(B\)，\(C\) 的对边分别是 \(a\)，\(b\)，\(c\)，已知 \(b=c\)，\(a^2=2b^2\left(1-\sin A\right)\)，则 \(A=\)（\quad）

\choices
    {\(\dfrac{3\pi}{4}\)}
    {\(\dfrac{\pi}{3}\)}
    {\(\dfrac{\pi}{4}\)}
    {\(\dfrac{\pi}{6}\)}
\end{problem}

\begin{problem}
已知函数 \(f(x)\) 的定义域为 \(\R\). 当 \(x<0\) 时，\(f(x)=x^3-1\)；当 \(-1\leq x\leq 1\) 时，\(f(-x)=-f(x)\)；当 \(x>\dfrac12\) 时，\(f\left(x+\dfrac12\right)=f\left(x-\dfrac12\right)\). 则 \(f(6)=\)（\quad）

\choices
    {\(-2\)}
    {\(-1\)}
    {\(0\)}
    {\(2\)}
\end{problem}

\begin{problem}
若函数 \(y=f(x)\) 的图象上存在两点，使得函数的图象在这两点处的切线互相垂直，则称 \(y=f(x)\) 具有 \(T\) 性质. 下列函数中具有 \(T\) 性质的是（\quad）

\choices
    {\(y=\sin x\)}
    {\(y=\ln x\)}
    {\(y=\e^x\)}
    {\(y=x^3\)}
\end{problem}

\section{填空题}

\begin{problem}
执行如图的程序框图，若输入 \(n\) 的值为 \(3\)，则输出的 \(S\) 的值为\fillinblank{}.

\bitmapfigure[max width=6.0cm]{img/2016/shandong_liberal/q11_fig1.png}
\end{problem}

\begin{problem}
观察下列等式：

\(\left(\sin\dfrac{\pi}{3}\right)^{-2}+\left(\sin\dfrac{2\pi}{3}\right)^{-2}=\dfrac43\times1\times2\)；

\(\left(\sin\dfrac{\pi}{5}\right)^{-2}+\left(\sin\dfrac{2\pi}{5}\right)^{-2}+\left(\sin\dfrac{3\pi}{5}\right)^{-2}+\left(\sin\dfrac{4\pi}{5}\right)^{-2}=\dfrac43\times2\times3\)；

\(\left(\sin\dfrac{\pi}{7}\right)^{-2}+\left(\sin\dfrac{2\pi}{7}\right)^{-2}+\left(\sin\dfrac{3\pi}{7}\right)^{-2}+\cdots+\left(\sin\dfrac{6\pi}{7}\right)^{-2}=\dfrac43\times3\times4\)；

\(\left(\sin\dfrac{\pi}{9}\right)^{-2}+\left(\sin\dfrac{2\pi}{9}\right)^{-2}+\left(\sin\dfrac{3\pi}{9}\right)^{-2}+\cdots+\left(\sin\dfrac{8\pi}{9}\right)^{-2}=\dfrac43\times4\times5\).

照此规律，
\(\left(\sin\dfrac{\pi}{2n+1}\right)^{-2}+\left(\sin\dfrac{2\pi}{2n+1}\right)^{-2}+\left(\sin\dfrac{3\pi}{2n+1}\right)^{-2}+\cdots+\left(\sin\dfrac{2n\pi}{2n+1}\right)^{-2}=\)\fillinblank{}.
\end{problem}

\begin{problem}
已知向量 \(\bs a=(1,-1)\)，\(\bs b=(6,-4)\). 若 \(\bs a\perp(t\bs a+\bs b)\)，则实数 \(t\) 的值为\fillinblank{}.
\end{problem}

\begin{problem}
已知双曲线 \(E:\dfrac{x^2}{a^2}-\dfrac{y^2}{b^2}=1\)（\(a>0\)，\(b>0\)）. 矩形 \(ABCD\) 的四个顶点在 \(E\) 上，\(AB\)，\(CD\) 的中点为 \(E\) 的两个焦点，且 \(2|AB|=3|BC|\)，则 \(E\) 的离心率是\fillinblank{}.
\end{problem}

\begin{problem}
已知函数
\(f(x)=\begin{cases}
|x|,&x\leq m,\\
x^2-2mx+4m,&x>m,
\end{cases}\)
其中 \(m>0\). 若存在实数 \(b\)，使得关于 \(x\) 的方程 \(f(x)=b\) 有三个不同的根，则 \(m\) 的取值范围是\fillinblank{}.
\end{problem}

\section{解答题}

\begin{problem}
某儿童乐园在“六一”儿童节推出了一项趣味活动. 参加活动的儿童需转动如图所示的转盘两次，每次转动后，待转盘停止转动时，记录指针所指区域中的数. 设两次记录的数分别为 \(x\)，\(y\). 奖励规则如下：

\begin{circlelist}
\item 若 \(xy\leq 3\)，则奖励玩具一个；
\item 若 \(xy\geq 8\)，则奖励水杯一个；
\item 其余情况奖励饮料一瓶.
\end{circlelist}

假设转盘质地均匀，四个区域划分均匀. 小亮准备参加此项活动.

\begin{enumerate}
\item 求小亮获得玩具的概率；
\item 请比较小亮获得水杯与获得饮料的概率的大小，并说明理由.
\end{enumerate}

\begingroup
\leftskip=0pt\relax
\bitmapfigure[max width=4.8cm]{img/2016/shandong_liberal/q16_fig1.png}
\endgroup
\end{problem}

\begin{problem}
设 \(f(x)=2\sqrt3\sin\left(\pi-x\right)\sin x-\left(\sin x-\cos x\right)^2\).

\begin{enumerate}
\item 求 \(f(x)\) 的单调递增区间；
\item 把 \(y=f(x)\) 的图象上所有点的横坐标伸长到原来的 \(2\) 倍（纵坐标不变），再把得到的图象向左平移 \(\dfrac{\pi}{3}\) 个单位，得到函数 \(y=g(x)\) 的图象，求 \(g\left(\dfrac{\pi}{6}\right)\) 的值.
\end{enumerate}
\end{problem}

\begin{problem}
在如图所示的几何体中，\(D\) 是 \(AC\) 的中点，\(EF\parallel DB\).

\begin{enumerate}
\item 已知 \(AB=BC\)，\(AE=EC\). 求证：\(AC\perp FB\)；
\item 已知 \(G\)，\(H\) 分别是 \(EC\) 和 \(FB\) 的中点. 求证：\(GH\parallel\) 平面 \(ABC\).
\end{enumerate}

\bitmapfigure[max width=4.8cm]{img/2016/shandong_liberal/q18_fig1.png}
\end{problem}

\begin{problem}
已知数列 \(\{a_n\}\) 的前 \(n\) 项和 \(S_n=3n^2+8n\)，\(\{b_n\}\) 是等差数列，且 \(a_n=b_n+b_{n+1}\).

\begin{enumerate}
\item 求数列 \(\{b_n\}\) 的通项公式；
\item 令 \(c_n=\dfrac{(a_n+1)^{n+1}}{(b_n+2)^n}\)，求数列 \(\{c_n\}\) 的前 \(n\) 项和 \(T_n\).
\end{enumerate}
\end{problem}

\begin{problem}
设函数 \(f(x)=x\ln x-ax^2+(2a-1)x\)，\(a\in\R\).

\begin{enumerate}
\item 令 \(g(x)=f'(x)\)，求函数 \(g(x)\) 的单调区间；
\item 已知 \(f(x)\) 在 \(x=1\) 处取得极大值，求实数 \(a\) 的取值范围.
\end{enumerate}
\end{problem}

\begin{problem}
已知椭圆 \(C:\dfrac{x^2}{a^2}+\dfrac{y^2}{b^2}=1\)（\(a>b>0\)）的长轴长为 \(4\)，焦距为 \(2\sqrt2\).

\begin{enumerate}
\item 求椭圆 \(C\) 的方程；
\item 过动点 \(M(0,m)\)（\(m>0\)）的直线交 \(x\) 轴于点 \(N\)，交 \(C\) 于点 \(A\)，\(P\)（\(P\) 在第一象限），且 \(M\) 是线段 \(PN\) 的中点. 过点 \(P\) 作 \(x\) 轴的垂线交 \(C\) 于另一点 \(Q\)，延长线 \(QM\) 交 \(C\) 于点 \(B\).
    \begin{enumerate}
    \item 设直线 \(PM\)，\(QM\) 的斜率分别为 \(k\)，\(k'\)，证明 \(\dfrac{k'}{k}\) 为定值；
    \item 求直线 \(AB\) 的斜率的最小值.
    \end{enumerate}
\end{enumerate}

\bitmapfigure[max width=4.8cm]{img/2016/shandong_liberal/q21_fig1.png}
\end{problem}
